\documentclass{article}

\usepackage{amsmath}
\usepackage{enumerate}
\usepackage[spanish]{babel}
\usepackage[utf8]{inputenc}
\usepackage[left=3cm,right=2.5cm,top=2.5cm,bottom=2.5cm]{geometry}

\begin{document}

\title{Práctica de Matemáticas Financiera II}
\author{César Andrés Cabrera Cabero}
\date{UPB}
\maketitle

\section{Interés Simple}

\begin{enumerate}
    \item Juan invirtió \$2000 a una tasa de interés simple del 8\% anual. ¿Cuánto interés ganará después de 3 años?
    
    \textbf{\$480}
    
    \item Calcula el tiempo necesario para que \$5000 a una tasa de interés simple del 6\% anual produzca \$800 de interés.
    
    \textbf{2 años y 8 meses}
    
    \item ¿Cuál es el interés simple ganado en un préstamo de \$3000 con una tasa de interés del 4\% anual durante 2 años?
    
    \textbf{\$240}
    
    \item María tomó un préstamo de \$4000 con una tasa de interés simple del 10\% anual. Si debe pagar \$480 de interés, ¿cuál es el tiempo del préstamo?
    
    \textbf{1 año con 2 meses y casi dos semanas}
    
    \item Si un préstamo de \$6000 con una tasa de interés del 5\% anual se pagó en su totalidad después de 2 años, ¿cuál es el monto total a pagar?
    
    \textbf{\$6600}
    
    \item Una persona invirtió \$800 a una tasa de interés simple del 12\% anual. Si ganó \$192 de interés, ¿cuánto tiempo mantuvo la inversión?
    
    \textbf{2 años}
    
    \item Si se depositan \$100 al final de cada mes en una cuenta de ahorro que paga un interés simple del 4\% anual, ¿cuánto se habrá acumulado después de 1 año?
    
    \textbf{\$2600}
    
    \item ¿Cuál será el interés simple ganado en un préstamo de \$2500 con una tasa de interés del 7\% anual durante 4 años?
    
    \textbf{\$700}
    
    \item Si se invierte \$3000 a una tasa de interés simple del 9\% anual, ¿cuánto tiempo tomará para que el monto se duplique?
    
    \textbf{8 años}
    
    \item Pedro invirtió \$1200 a una tasa de interés simple del 3\% anual. Si ganó \$108 de interés, ¿cuánto tiempo mantuvo la inversión?
    
    \textbf{1 año}
\end{enumerate}

\section{Interés Compuesto}

\begin{enumerate}
    \item Calcula el monto final después de 5 años si se invierten \$5000 a una tasa de interés compuesto del 6\% anual.
    
    \textbf{\$6730.73}
    
    \item ¿Cuánto tiempo tomará para que una inversión de \$8000 a una tasa de interés compuesto del 9\% anual duplique su valor?
    
    \textbf{Aproximadamente 8 años}
    
    \item Si se depositan \$200 al final de cada mes en una cuenta de ahorros que paga un interés compuesto del 5\% anual, ¿cuánto se habrá acumulado después de 2 años?
    
    \textbf{\$6385.55}
    
    \item María depositó \$1000 al principio de cada año en una cuenta de ahorros que paga un interés compuesto del 8\% anual. ¿Cuál será el saldo de la cuenta después de 4 años?
    
    \textbf{\$1271.94}
    
    \item Calcula el monto total acumulado si se invierten \$3000 a una tasa de interés compuesto del 7\% anual durante 3 años.
    
    \textbf{\$3726.10}
    
    \item Si se invierte \$400 a una tasa de interés compuesto del 10\% anual, ¿cuál será el monto acumulado después de 4 años?
    
    \textbf{\$484.00}
    
    \item Si una inversión inicial de \$6000 crece a \$9000 en 6 años, ¿cuál es la tasa de interés compuesto anual?
    
    \textbf{5\%}
    
    \item Si una inversión inicial de \$1000 crece a \$1380 en 2 años, ¿cuál es la tasa de interés compuesto anual?
    
    \textbf{9\%}
    
    \item Calcula el monto final después de 8 años si se invierten \$2000 a una tasa de interés compuesto del 5\% anual.
    
    \textbf{\$2973.64}
    
    \item Si se invierte \$1200 a una tasa de interés compuesto del 4\% anual, ¿cuánto tiempo tomará para que el monto se duplique?
    
    \textbf{18 años}
\end{enumerate}
\section{Gradientes}
\begin{enumerate}
    \item Considere una anualidad anticipada con gradiente aritmético cuyo valor inicial es de Bs. 450 y el incremento es de Bs 25 cada mes, además se sabe que la cuenta de ahorro paga un interés del 8\% anual capitalizable trimestralmente. ¿Cuánto se habrá ahorrado después de 6.75 años?
    \item Considere una anualidad anticipada con gradiente geométrico cuyo valor inicial es de Bs 120 y la razón de cambio es del 20\% (incremento), además se sabe que la cuenta de ahorro paga un interes del 6.5\% anual capitalizable cuatrimestralmente. ¿Cuanto se habrá ahorrado después de 3.25 años?
    \item Sergio quiere empezar a ahorrar en un banco que paga el 5\% semestral capitalizable bimestralmente. Él cuenta con Bs. 650 cada mes y quiere ahorrar durante 25 años. Despúes los 25 años quiere retirar una anualidad con gradiente geométrico, la anualidad empezará con Bs 1 000 con razón del 95\% mensual, ¿cuánto tiempo podrá retirar esa anualidad con gradiente hasta que se quede sin fondos?
    \item Consideremos una anualidad anticipada con gradiente geométrico que empieza con Bs 540 y razon de cambio del \%10 mensual, ésta se invierte en una caja de ahorro que paga el 4\% semestral efectivo durante 5 años. Si la anualidad con gradiente sería aritmético y empezaría con Bs 540, ¿cuánto tendría que ser el incremento para que, en el mismo tiempo y con la misma tasa de interés, se alcance el mismo monto final?  
\end{enumerate}
\section{Anualidades}

\begin{enumerate}
    \item Calcula el valor futuro de una anualidad anticipada de \$2 000 cada año durante 5 años con una tasa de interés del 7\% efectivo anual.
    
    % \textbf{\$11026.34}
    
    \item ¿Cuánto debe depositar Juan al final de cada trimestre en una cuenta que paga un interés del 8\% anual capitalizable mensualmente para tener \$10 000 después de 3 años?
    
    % \textbf{\$197.47}
    
    \item María planea jubilarse en 25 años y quiere tener \$500 000 ahorrados. Si puede ahorrar \$200 al mes en una cuenta que paga un interés del 6\% anual capitalizable trimestralmente, ¿podrá alcanzar su objetivo?
    
    % \textbf{No, solo alcanzaría \$144,442.46}
    
    \item Si se depositan \$150 al principio de cada mes en una cuenta de ahorro que paga un interés del 5\% anual capitalizable bimestralmente, ¿cuánto se habrá acumulado después de 10 años?
    
    % \textbf{\$20,643.57}
    
    \item Tenemos un saldo $k$ dólares en una cuenta, si decidimos retirar un monto de \$3 000 al final de cada año durante 7 años con una tasa de interés del 8\% anual capitalizable cuatrimestralmente, y así vaciar la cuenta, ¿cuánto había en un principio?
    
    \item Contamos con un monto de \$10 000 en una caja de ahorro que paga el 3.5\% anual capitalizable mensualmente de la cual vamos retirando cada semestre vencido de \$580 durante 6.5 años, depués decidimos retirar \$230 cada mes vencido durante 2.25 años, ¿cuánto dinero nos queda al final los 8.75 años?
    \item Si contamos con \$$k$ dólares en una caja de ahorro que paga \%5.5 semestral efectivo y queremos sacar de dicha cuenta $k/50$ dólares mensuales, ¿cuánto tiempo después la cuenta tendrá la mitad del saldo inicial?
    
    % \textbf{\$18926.83}
    
\end{enumerate}
\section{Amortizaciones}
    \begin{enumerate}
        \item Alejandra está interesada en un crédito de Bs. 1 450 530 y el Banco Chinpilin le ofrece una tasa de interés del 4.45\% anual capitalizable mensualmente. Considerando que Alejandra solo cuenta con Bs. 18 850 de salario y el informe de riesgo establece que solo se puede usar hasta el 40\% de su salario como capacidad de pago, se plantean las siguientes interrogantes:
        \begin{enumerate}
            \item ¿Cuántos meses tiene que pagar su deuda y cuál sería su cuota mensual si decide aplicar a un plan de cuota fija?\label{a}
            \item Si Alejandra quisiera pagar en la mitad del tiempo calculado en \ref{a} y manteniendo el plan de cuotas fija, ¿cuánto tendría que aumentar su cuota?
            \item Alejandra sabe que tendrá un incremeto salarial del 8\% cada cinco años, 60 meses, y quiere saber que le conviene más, si un plan de cuota fija o un plan de cuota variable(considerando como primer pago su capacidad de pago).
            \item Establezca el plan de pagos de Alejandra en ambos planes de pago si la tasa de interés del banco aumentará después de tres años a 5.5\% anual capitalizable mensualmente.
            \item Realice un gráfico comparativo con todo lo anterior.
        \end{enumerate}
    \end{enumerate}

\section{Evaluación de proyectos}
\begin{enumerate}
    \item Fuji Software, Inc., tiene los siguientes proyectos mutuamente excluyentes:

    \begin{center}
    \begin{tabular}{|c|c|c|}
        \hline
        Año & Proyecto A & Proyecto B \\
        \hline
        0 & $-10,000$ & $-12,000$ \\
        1 & $3,500$ & $3,000$ \\
        2 & $4,000$ & $2,000$ \\
        3 & $1,800$ & $3,000$ \\
        4 & $2,500$ & $3,500$ \\
        5 & $1,000$ & $4,000$ \\
        6 & $1,800$ & $5,000$ \\
        \hline
    \end{tabular}
    \end{center}
    
    \begin{enumerate}[(a)]
        \item Suponga que el límite del periodo de recuperación de Fuji es de 4 años y la tasa de descuento es del 8.5\%. ¿Cuál de estos dos proyectos debería elegir la empresa?
        \item Suponga que Fuji usa la regla del valor presente neto para clasificar estos dos proyectos. ¿Qué proyecto debe elegir la empresa si la tasa de descuento apropiada es del 9\%?
    \end{enumerate}
    
    \item Un proyecto de inversión proporciona flujos de ingreso efectivo de \$70 dólares por año durante ochenta años con una tasa descuento del 0.15\%. ¿Cuál es el periodo de recuperación del proyecto si el costo inicial es de \$4,100 dólares? ¿Y si el costo inicial es de \$6,200 dólares? ¿Y si es de \$8,000 dólares?
    
    \item Un proyecto de inversión tiene flujos anuales de ingreso efectivo de \$6,000, \$6,500, \$7,000 y \$8,000 dólares, y una tasa de descuento del 9\%. ¿Cuál es el periodo de recuperación descontado de estos flujos efectivos si el costo inicial es de \$8,000 dólares?
    \item Un proyecto tiene un costo inicial de $I$, un rendimiento requerido de $R$ y paga $C$ anualmente durante $N$ años.

    \begin{enumerate}[(a)]
        \item Encuentre $C$ en términos de $I$, $N$ y $R$ de tal modo que el proyecto tenga un periodo de recuperación justo igual a su vida y realice un gráfico que ilustre lo anterior.
        \item Encuentre $C$ en términos de $I$, $N$ y $R$ de tal modo que este sea un proyecto rentable de acuerdo con la regla de decisión del VPN y realice un gráfico que ilustre lo anterior.
        \item Encuentre $C$ en términos de $I$, $N$ y $R$ de tal modo que el proyecto tenga una razón del costo-beneficio de 2y realice un gráfico que ilustre lo anterior.
    \end{enumerate}
    
    \item Suponga que hoy le ofrecen \$8,000 dólares, pero debe hacer los siguientes pagos:
    
    \begin{center}
    \begin{tabular}{|c|c|}
        \hline
        Año & Flujos de efectivo (\$) \\
        \hline
        0 & 8000 \\
        1 & -4400 \\
        2 & -2700 \\
        3 & -1900 \\
        4 & -1500 \\
        \hline
    \end{tabular}
    \end{center}
    
    \begin{enumerate}[(a)]
        \item ¿Cuál es la TIR de esta oferta?
        \item Si la tasa de descuento apropiada es del 10\%, ¿debe aceptar esta oferta?
        \item Si la tasa de descuento apropiada es del 20\%, ¿debe aceptar esta oferta?
    \end{enumerate}
\end{enumerate}

\section{Bonos}
\begin{enumerate}
    \item Usted acaba de adquirir una obligación del Estado, que paga un 6\% de interés anual y a la que le quedan cinco años de vida, por un valor de Bs 95. Sabiendo que el valor nominal es de Bs 100. Construya una tabla de pagos para los cinco años restantes.
    \item Usted invierte en un par de bonos cupón-cero. Uno de ellos madura dentro de un año pagando 100 euros. Su precio actual es de 86,93 euros. El otro madura dentro de dos años pagando 110 euros, y su precio actual es de 94,307 euros. Cálcular el rendimiento de cada Bonos.
    
    \textbf{Respuesta:} $r_1=15.035\%$ y $r_2=8\%$
    \item Suponga que una compañía alemana emite un bono con valor nominal de 1.000 euros, 20 años para el vencimiento y tasa de cupón de 8.4\% pagadero anualmente. Si el rendimiento al vencimiento es de 15\%, ¿cuál es el precio actual del bono?
    \item Calcule el rendimiento al vencimiento de un bono que actualmente cotiza en el mercado a \$750,75579 con las siguientes características: Valor Nominal: \$1.000, Tiempo al vencimiento: 10 años, Tasa cupón de 6\%, Pago de cupones semestrales.    
\end{enumerate}
\end{document}
